\origin{ai}
\subsection{Finite-window data required for a theta-square factorization certificate}
\label{sec:window-fibonacci-cube-theta-square-factorization-certificate}

\paragraph{Available datum.}
The present record names the intended object as a finite forced-window construction on a Fibonacci-cube carrier, with a theta-square factorization certificate as the intended output. No finite row data are supplied with the record. In particular, the record does not specify a word length, a forbidden adjacency rule, a list of admissible words, a partition of those words, a finite relation on the carrier, or a matrix whose square, compression, trace, determinant, rank, kernel, or spectrum can be checked directly. Therefore this text can only record the exact obstruction to forming the certificate, not the certificate itself.

\paragraph{Carrier data that must be fixed.}
A self-contained finite certificate must first give a concrete carrier such as a displayed subset $X \subseteq \{0,1\}^n$, with $n$ explicitly stated and with the membership predicate written from first principles. For a Fibonacci-cube window this would normally mean a rule of the form: a word $x=(x_1,\ldots,x_n)$ belongs to $X$ exactly when each $x_i$ is either $0$ or $1$ and a specified local exclusion holds for adjacent positions. The present record does not state $n$ and does not state the local exclusion. Consequently neither $|X|$ nor any ordered list of elements of $X$ is determined by the supplied facts.

\paragraph{Partition and relation data that must be fixed.}
The advertised factorization also requires finite structure on the carrier. At minimum one must know whether the factorization is taken over all of $X$, over a partition $X=P_1\sqcup\cdots\sqcup P_r$, or over a quotient by an explicitly listed equivalence relation. A finite relation $R\subseteq X\times X$ must also be specified if the certificate is to use adjacency, compatibility, transition, or incidence data. Without the blocks $P_i$ and without the ordered pairs in $R$, no incidence matrix can be reconstructed and no factorization identity has a well-defined left-hand side.

\paragraph{Matrix data that must be fixed.}
A theta-square certificate must specify the finite matrix or kernel being squared. This means giving entries $A_{xy}$ for every ordered pair $(x,y)\in X\times X$, or giving an equivalent rule that uniquely computes those entries. The asserted identity must then be a concrete equality such as
$$
\begin{aligned}
B_{xz} &= \sum_{y\in X} A_{xy}A_{yz}
\end{aligned}
$$
for all $x,z\in X$, together with the exact target matrix $B$. If the certificate is spectral, the exact characteristic polynomial, eigenvalue multiplicities, kernel basis, or block factorization must be stated as finite data. None of these rows, entries, or target identities are present in the supplied facts.

\paragraph{Witness extraction.}
A completed certificate must say how a reader verifies the construction by finite inspection: list the carrier, list the relation or matrix entries, perform the stated finite sums, and compare every resulting entry with the target table. The present record does not include the tables needed for that extraction. Hence no exact count, no exact matrix identity, no square-root relation, and no factorization witness can be honestly asserted here.

\paragraph{Local boundary.}
This record does not establish any physical-constant identification, global flow law, external dependency, metrological comparison, or interpretation beyond a finite construction. Even a numerical agreement, if later supplied, would remain separate from the finite algebraic certificate unless an additional certificate states the comparison protocol and the permitted tolerance.